\chapter{Calculation of the collision kernel}

\section{Linearization of the collision integral}

The Boltzmann equation is
\begin{equation}
  \odv{f}{t} + \vv{v} \vdot \odv{f}{\vv{r}} + \vv{F} \vdot \odv{f}{\hbar\vv{p}} = \mathcal{J}[f]
\end{equation}
where \(f(t, \vv{x}, \vv{p})\) is the one particle density function, and \(\mathcal{J}\), the collision integral given by
\begin{equation}
  \mathcal{J}[f]
  = - \frac{1}{(2\cpi)^6} \idotsint \odif{\vv{p}_2,\vv{p_1}',\vv{p_2}'} W(\vv{p}_1', \vv{p}_2' | \vv{p}_1, \vv{p}_2) \mathcal{B}
\end{equation}
where
\begin{multline}
  \mathcal{B} = \big[ f(\vv{p}_1)f(\vv{p}_2)(1 - f(\vv{p}_1'))(1 - f(\vv{p}_2')) \\
    - f(\vv{p}_1')f(\vv{p}_2')(1 - f(\vv{p}_1))(1 - f(\vv{p}_2)) \big]
\end{multline}
This integral is non-linear and is hard to solve. Hence in this section, we shall simplify the integral by linearizing it w.r.t. \(f\).

From here, we shall write \(f_i = f(\vv{p}_i)\) and \(f_{0,i} = f(\vv{p}_i)\). Primed \(f_i'\) are understood as \(f(\vv{p}_i')\).

Consider the taylor expansion of \(\mathcal{B}\) at \(f_0\):
\begin{equation}
  \mathcal{B} = \mathcal{B}|_{f_0} + \mathcal{B}_{\textrm{linear}} + \mathcal{B}_{\textrm{non-linear}}.
\end{equation}
The linear part can be expanded w.r.t. each of the \(f_i\)'s to get
\begin{equation}
  \mathcal{B}_{\textrm{linear}}
  = \fdif{f_1}\pdv{\mathcal{B}}{f_1}_{f_0} + \fdif{f_2}\pdv{\mathcal{B}}{f_2}_{f_0}
  + \fdif{f_1'}\pdv{\mathcal{B}}{f_1'}_{f_0} + \fdif{f_2'} \pdv{\mathcal{B}}{f_2'}_{f_0}
\end{equation}

% \begin{notation}{Notation}{}
%   \vspace{-1.5em}
%   \begin{align}
%     \vv{S} & = \vv{p}_1 + \vv{p}_2                         \\
%     \vv{D} & = \vv{p}_1 - \vv{p}_2                         \\
%     \vv{q} & = \vv{p}_1' - \vv{p}_1 = \vv{p}_2 - \vv{p}_2' \\
%     R      & = \frac{1}{4}(D_y^2 - D_x^2)
%   \end{align}
% \end{notation}
%
% \section{Kinematic constraints}
%
% These kinematic constraints includes the momentum and energy conservation which will be imposed into the integral via the use of dirac deltas:
% \begin{gather}
%   \vv{p}_1 + \vv{p}_2 - \vv{p}_1' - \vv{p}_2' = 0 \\
%   \epsilon(\vv{p}_1) + \epsilon(\vv{p}_2) - \epsilon(\vv{p}_1') - \epsilon(\vv{p}_2') = 0
% \end{gather}
% The momentum constraint is trivial, and can be satisfied by setting
% \begin{gather}
%   \vv{p}_1' = \vv{p}_1 + \vv{q}, \\
%   \vv{p}_2' = \vv{p}_2 - \vv{q}.
% \end{gather}
% Hence, we move to the energy constraint with the dispersion relation given by
% \begin{equation}
%   \epsilon(\vv{p}) = \frac{\hbar^2}{2m}(p_x^2 - p_y^2).
% \end{equation}
% where \(m\) is the effective mass, carrying the units of \(\unit{\kilogram^3\meter^3\second^{-2}}\).
%
% The expression of the energy constraint can be expanded to
% \begin{align}
%    & \begin{multlined}[t][0.8\textwidth]
%        \frac{\hbar^2}{2m} \Big[
%          (p_{1x}^2 - p_{1y}^2) + (p_{2x}^2 - p_{2y}^2) \\
%          - \ab((p_{1x} + q_x)^2 - (p_{1y} + q_y)^2) \\
%          - \ab((p_{2x} - q_x)^2 - (p_{2y} - q_y)^2)
%          \Big]
%      \end{multlined}                               \\
%    & = \begin{multlined}[t][0.8\textwidth]
%          \frac{\hbar^2}{2m} \Big[
%            p_{1x}^2 - p_{1y}^2 + p_{2x}^2 - p_{2y}^2 \\
%            - p_{1x}^2 - 2p_{1x}q_x - q_x^2 + p_{1y}^2 + 2p_{1y}q_y + q_y^2 \\
%            -p_{2x}^2 + 2p_{2x}q_x - q_x^2 + p_{2y}^2 - 2p_{2y}q_y + q_y^2
%            \Big]
%        \end{multlined} \\
%    & = \frac{\hbar^2}{2m}[-2p_{1x}q_x + 2p_{2x}q_x - 2q_x^2 + 2p_{1y}q_y - 2p_{2y}q_y + 2q_y^2]                                                                              \\
%    & = \frac{\hbar^2}{m}\ab\Big[\ab\big(q_y^2 + q_y(p_{1y} - p_{2y})) - \ab\big(q_x^2 + q_x(p_{1x} - p_{2x}))]
%   \intertext{Let \(\vv{D} = \vv{p}_1 - \vv{p}_2\), and complete the square:}
%    & = \frac{\hbar^2}{m}\ab[q_y(D_y + q_y) - q_x(D_x + q_x)]                                                                                                                 \\
%    & = \begin{multlined}[t]
%          \frac{\hbar^2}{m}\Bigg[\ab(q_y^2 + 2(q_y)\ab\Big(\frac{1}{2}D_y) + \frac{1}{4}D_y^2)
%            - \ab(q_x^2 + 2(q_x)\ab\Big(\frac{1}{2}D_x) + \frac{1}{4}D_x^2) \\
%            - \frac{1}{4}(D_y^2 - D_x^2)\Bigg]
%        \end{multlined}                                                                      \\
%    & = \frac{\hbar^2}{m} \ab[\ab(q_y + \frac{D_y}{2})^2 - \ab(q_x + \frac{D_x}{2})^2 - \frac{1}{4}(D_y^2 - D_x^2)]
% \end{align}
%
% Here, we can define
% \begin{equation}
%   R = \frac{1}{4}(D_y^2 - D_x^2)
% \end{equation}
%
% % We then perform a change of coordinates in the scheme of
% % \begin{gather}
% %   r\cosh\theta = q_y + \frac{D_y}{2} \\
% %   r\sinh\theta = q_x + \frac{D_x}{2} \\
% %   R = \frac{1}{4}(D_y^2 - D_x^2)
% % \end{gather}
% To evaluate the delta functions, we split the phase space into four cones that are analogous to the light cone structure:
% \begin{align}
%   \begin{cases}
%     \textrm{Future:} & R > 0, D_y > 0 \\
%     \textrm{Past:}   & R > 0, D_y < 0 \\
%     \textrm{Right:}  & R < 0, D_x > 0 \\
%     \textrm{Left:}   & R < 0, D_x < 0 \\
%   \end{cases}
% \end{align}
%
% For the future and past cone (\(R > 0\)), we set
% \begin{gather}
%   r\cosh\theta = q_y + \frac{D_y}{2} \\
%   r\sinh\theta = q_x + \frac{D_x}{2}
% \end{gather}
% Combining the two equations prior, we get
% \begin{align}
%   r^2 (\cosh^2 \theta - \sinh^2 \theta) & = \ab(q_y + \frac{D_y}{2})^2 + \ab(q_x + \frac{D_x}{2})^2             \\
%   r                                     & = \pm \sqrt{\ab(q_y + \frac{D_y}{2})^2 + \ab(q_x + \frac{D_x}{2})^2},
% \end{align}
% where \(r, \theta \in \mathbb{R}\). As seen, the integral is split into the \(r > 0\) branch and \(r < 0\) branch. Either ways, we get
% \begin{equation}
%   \frac{\hbar^2}{m}\ab[r^2 (\cosh^2\theta - \sinh^2\theta) - R] = \frac{\hbar^2}{m}(r^2 - R)
% \end{equation}
% Substitute this into the collision integral to get:
% \begin{equation}
%   \mathcal{J}[f] = \begin{multlined}[t][0.8\textwidth]
%     - \frac{1}{(2\cpi)^6} \idotsint \odif{\vv{p}_2,\vv{p}_1',\vv{p}_2'} \\
%     \times \frac{(2\cpi)^3}{\hbar} |V^2| \ddel(\vv{p}_1 + \vv{p}_2 - \vv{p}_1' - \vv{p}_2') \frac{m}{\hbar^2} \ddel(r^2 - R) \\
%     \times F_{121'2'}\ab\big[\psi(\vv{p}_1) + \psi(\vv{p}_2) - \psi(\vv{p}_1') - \psi(\vv{p}_2')].
%   \end{multlined}
% \end{equation}
% For simplicity, let \(\mathcal{C} = F_{121'2'}\ab\big[\psi(\vv{p}_1) + \psi(\vv{p}_2) - \psi(\vv{p}_1') - \psi(\vv{p}_2')]\), the integral then simplifies to
% \begin{equation}
%   - \frac{m|V|^2}{(2\cpi)^3\hbar^3} \idotsint \odif{\vv{p}_2,\vv{p}_1',\vv{p}_2'} \ddel(\vv{p}_1 + \vv{p}_2 - \vv{p}_1' - \vv{p}_2') \ddel(r^2 - R) \mathcal{C}
% \end{equation}
%
% To evaluate these delta functions, we perform another change of coordinates from \((\vv{p}_1', \vv{p}_2')\) to \((\vv{S}', \vv{q})\) given by
% \begin{gather}
%   S_x' = p_{1x}' + p_{2x}' \\
%   S_y' = p_{1y}' + p_{2y}' \\
%   q_x = p_{1x}' - p_{1x} \\
%   q_y = p_{1y}' - p_{1y}
% \end{gather}
% The jacobian of this transformation is then
% \begin{equation*}
%   \ab|\jdv{S_x', S_y', q_x, q_y}{p_{1x}', p_{1y}', p_{2x}', p_{2y}'}| = \begin{vmatrix}
%     \pdv{S_x'}{p_{1x}'} & \pdv{S_x'}{p_{1y}'} & \pdv{S_x'}{p_{2x}'} & \pdv{S_x'}{p_{2y}'} \\
%     \pdv{S_y'}{p_{1x}'} & \pdv{S_y'}{p_{1y}'} & \pdv{S_y'}{p_{2x}'} & \pdv{S_y'}{p_{2y}'} \\
%     \pdv{q_x}{p_{1x}'}  & \pdv{q_x}{p_{1y}'}  & \pdv{q_x}{p_{2x}'}  & \pdv{q_x}{p_{2y}'}  \\
%     \pdv{q_y}{p_{1x}'}  & \pdv{q_y}{p_{1y}'}  & \pdv{q_y}{p_{2x}'}  & \pdv{q_y}{p_{2y}'}  \\
%   \end{vmatrix} = \begin{vmatrix}
%     1 & 0 & 1 & 0 \\
%     0 & 1 & 0 & 1 \\
%     1 & 0 & 0 & 0 \\
%     0 & 1 & 0 & 0
%   \end{vmatrix} = 1
% \end{equation*}
% Hence, \(\odif{\vv{p}_1',\vv{p}_2'} = \odif{\vv{S}',\vv{q}'}\). Substituting this into the integral to get:
% \begin{align}
%   \mathcal{J}[f] & = -\frac{m|V|^2}{(2\cpi)^3\hbar^3} \idotsint \odif{\vv{p}_2,\vv{S}',\vv{q}} \ddel(\vv{p}_1 + \vv{p}_2 - \vv{S}')\ddel(r^2 - R)\mathcal{C} \\
%                  & = -\frac{m|V|^2}{(2\cpi)^3\hbar^3} \idotsint \odif{\vv{p}_2,\vv{q}} \ddel(r^2 - R) \mathcal{C}|_{\vv{p}_1 + \vv{p}_2 = \vv{S}'}
% \end{align}
%
% The energy conservation delta function can be splitted in to by the identity
% \begin{equation}
%   \ddel(g(x)) = \sum_{i} \frac{\ddel(x - x_i)}{|g'(x_i)|}
% \end{equation}
% where \(x_i\) is the roots of \(g(x)\). Here, \(r^2 - R\) has roots \(\sqrt{R}\) and \(-\sqrt{R}\). Thus,
% \begin{align}
%   \ddel(r^2 - R) & = \frac{\ddel(r + \sqrt{R})}{2\sqrt{R}} + \frac{\ddel(r - \sqrt{R})}{2\sqrt{R}} \\
%                  & = \frac{1}{2\sqrt{R}}\ab(\ddel(r + \sqrt{R}) + \ddel(r - \sqrt{R}))
% \end{align}
% But since \(r \in \ocintv{0, \infty}\), only \(\ddel(r - \sqrt{R})\) contributes to the integral. Hence,
% \begin{equation}
%   \mathcal{J}[f] = - \frac{m|V|^2}{16 \cpi^3 \hbar^3} \idotsint \odif{\vv{p}_2,\vv{q}} \frac{1}{\sqrt{R}} \ddel(r - \sqrt{R}) \mathcal{C}|_{\vv{p}_1 + \vv{p}_2 = \vv{S}'}.
% \end{equation}
%
% To evaluate the final delta function, we transform from \((\vv{p}_2, \vv{q})\) to \((r, \theta, \vv{D})\) via
% \begin{gather}
%   r \cosh\theta = q_y + \frac{D_y}{2}, \label{eq:rcoshtheta-to-qy} \\
%   r \sinh\theta = q_x + \frac{D_x}{2}, \label{eq:rsinhtheta-to-qx} \\
%   D_x = p_{1x} - p_{2x}, \\
%   D_y = p_{1y} - p_{2y}.
% \end{gather}
% Since \(\vv{p}_1\) is fixed, \(\odif{\vv{D}} = \odif{\vv{p}_2}\). The only work here is then to calculate the Jacobian of the transformation \((\vv{q})\) to \((r, \theta)\).
% \begin{equation}
%   \ab|\jdv{r, \theta}{q_x, q_y}| = \begin{vmatrix}
%     \pdv{r}{q_x}      & \pdv{r}{q_y}      \\
%     \pdv{\theta}{q_x} & \pdv{\theta}{q_y} \\
%   \end{vmatrix} = \ab|\pdv{r}{q_x}\pdv{\theta}{q_y} - \pdv{r}{q_y}\pdv{\theta}{q_x}|
% \end{equation}
%
% To evaluate the partial derivative of \(r\) w.r.t. \(q_x\) and \(q_y\), we take the sum of squares of \cref{eq:rcoshtheta-to-qy} and \cref{eq:rsinhtheta-to-qx} to yield
% \begin{align}
%   r^2(\cosh^2\theta - \sinh^2\theta) = \ab(q_y + \frac{D_y}{2})^2 - \ab(q_x + \frac{D_x}{2})^2 \\
%   r^2 = \ab(q_y + \frac{D_y}{2})^2 - \ab(q_x + \frac{D_x}{2})^2,
% \end{align}
% in which
% \begin{align}
%   \pdv{r^2}{q_x}  & = \pdv*{\ab(\ab\Big(q_y + \frac{D_y}{2})^2 - \ab\Big(q_x + \frac{D_x}{2})^2)}{q_x} \\
%   2r \pdv{r}{q_y} & = - 2\ab(q_x + \frac{D_x}{2})                                                      \\
%   \pdv{r}{q_x}    & = -\frac{1}{r}\ab(q_x + \frac{D_x}{2}) = - \sinh\theta.
% \end{align}
% Similarly,
% \begin{align}
%   \pdv{r^2}{q_y}  & = \pdv*{\ab(\ab\Big(q_y + \frac{D_y}{2})^2 - \ab\Big(q_x + \frac{D_x}{2})^2)}{q_y} \\
%   2r \pdv{r}{q_y} & = 2\ab(q_y + \frac{D_y}{2})                                                        \\
%   \pdv{r}{q_y}    & = \frac{1}{r}\ab(q_y + \frac{D_y}{2}) = \cosh\theta.
% \end{align}
%
% To evaluate the partial derivative of \(\theta\) w.r.t. \(q_x\) and \(q_y\), take the quotient of \cref{eq:rsinhtheta-to-qx,eq:rcoshtheta-to-qy} to get
% \begin{equation}
%   \tanh \theta = \frac{q_x + D_x/2}{q_y + D_y/2},
% \end{equation}
% in which
% \begin{align}
%   \pdv{\tanh\theta}{q_x}          & = \pdv*{\ab(\frac{q_x + D_x/2}{q_y + D_y/2})}{q_x}             \\
%   \sech^2\theta \pdv{\theta}{q_x} & = \frac{1}{q_y + D_y/2}                                        \\
%   \pdv{\theta}{q_x}               & = \frac{\cosh^2\theta}{r\cosh\theta} = \frac{1}{r}\cosh\theta.
% \end{align}
% Similarly,
% \begin{align}
%   \pdv{\tanh\theta}{q_y}          & = \pdv*{\ab(\frac{q_x + D_x/2}{q_y + D_y/2})}{q_y}                               \\
%   \sech^2\theta \pdv{\theta}{q_y} & = - \frac{q_x + D_x/2}{(q_y + D_y/2)^2}                                          \\
%   \pdv{\theta}{q_y}               & = -\cosh^2\theta \frac{r\sinh\theta}{r^2\cosh^2\theta} = -\frac{1}{r}\sinh\theta
% \end{align}
%
% The Jacobian then evaluates to
% \begin{align}
%    & \ab|(-\sinh\theta) \ab(-\frac{1}{r}\sinh\theta) - (\cosh\theta) \ab(\frac{1}{r}\cosh\theta)| \\
%    & = \ab|\frac{1}{r}\ab(\sinh^2\theta - \cosh^2\theta)| = \frac{1}{r}
% \end{align}
% Hence, \(\odif{\vv{q}} = r\odif{r,\theta}\) and the integral simplifies to
% \begin{align}
%   \mathcal{J}[f] & = - \frac{m|V^2|}{16\cpi^3\hbar^3} \idotsint \odif{\vv{D},r,\theta} \cdot \frac{r}{\sqrt{R}}\ddel(r - \sqrt{R}) \mathcal{C}|_{\vv{p}_1 + \vv{p}_2 = \vv{S}'} \\
%                  & = - \frac{m|V^2|}{16\cpi^3\hbar^3} \idotsint \odif{D_x,D_y,\theta}~\mathcal{C}|_{\vv{p}_1 + \vv{p}_2 = \vv{S}', r = \sqrt{R}}.
% \end{align}

\section{Preliminary simplification of the Fermi factor}

Let \(a \coloneq \beta(\epsilon(\vv{p}_1) - \mu)\) and \(b \coloneq \beta(\epsilon(\vv{p}_2) - \mu)\), then
\begin{align}
  f_0(\vv{p}_1)f_0(\vv{p}_2)
   & = \frac{1}{(1 + \e^a)(1 + \e^b)}                                                      \\
   & = \frac{1}{1 + \e^a + \e^b + \e^{a + b}}                                              \\
   & = \frac{1}{
    \e^{\frac{a + b}{2}}\ab(\e^{-\frac{a + b}{2}} + \e^{\frac{a + b}{2}}
    + \e^{\frac{a - b}{2}} + \e^{-\frac{a - b}{2}})
  }                                                                                        \\
   & = \frac{1}{2\e^{\frac{a + b}{2}}
    \ab[\frac{1}{2}(\e^{-\frac{a + b}{2}} + \e^{\frac{a + b}{2}})
      +\frac{1}{2}\ab(\e^{\frac{a - b}{2}} + \e^{-\frac{a - b}{2}})]
  }                                                                                        \\
   & = \frac{1}{2\e^{\frac{a + b}{2}}\ab[\cosh(\frac{a + b}{2}) + \cosh(\frac{a - b}{2})]}
\end{align}
Similarly, let \(a' \coloneq \beta(\epsilon(\vv{p}_1') - \mu)\) and \(b' \coloneq \beta(\epsilon(\vv{p}_2') - \mu)\), then
\begin{align}
  \ab(1 - f_0(\vv{p}_1'))\ab(1 - f_0(\vv{p}_2'))
   & = \ab(\frac{\e^{a'}}{1 + \e^{a'}})\ab(\frac{\e^{b'}}{1 + \e^{b'}})                                     \\
   & = \frac{\e^{a' + b'}}{(1 + \e^{a'})(1 + \e^{b'})}                                                      \\
   & = \frac{\e^{a' + b'}}{2\e^{\frac{a' + b'}{2}}\ab[\cosh(\frac{a' + b'}{2}) + \cosh(\frac{a' - b'}{2})]}
\end{align}
Together, we get that
\begin{equation}
  F_{121'2'}
  = \frac{1}{2\e^{\frac{a + b}{2}} \ab[\cosh(\frac{a + b}{2}) + \cosh(\frac{a - b}{2})]}
  \cdot \frac{\e^{a' + b'}}{2\e^{\frac{a' + b'}{2}}\ab[\cosh(\frac{a' + b'}{2}) + \cosh(\frac{a' - b'}{2})]} \nonumber \\
\end{equation}
By energy conservation, \(a' + b' = a + b\). Hence, \(F_{121'2'}\) evaluates to
\begin{equation}
  \frac{1}{
    4\ab[\cosh \xi + \cosh\ab(\frac{\beta}{2}(\epsilon(\vv{p_1}) - \epsilon(\vv{p}_2)))]
    \ab[\cosh \xi + \cosh\ab(\frac{\beta}{2}(\epsilon(\vv{p}_1') - \epsilon(\vv{p}_2')))]
  }
\end{equation}
where \(\xi \coloneq \frac{\beta}{2}(\epsilon(\vv{p}_1) + \epsilon(\vv{p}_2) - 2\mu)\).

\section{Kinematic constraints}

The kinematic constraints includes the momentum and energy conservation which will be imposed into the integral via dirac deltas:
\begin{gather}
  \vv{p}_1 + \vv{p}_2 - \vv{p}_1' - \vv{p}_2' = 0, \\
  \epsilon(\vv{p}_1) + \epsilon(\vv{p}_2) - \epsilon(\vv{p}_1') - \epsilon(\vv{p}_2') = 0.
\end{gather}

To evaluate these dirac deltas, we turn into the center of momentum and relative momentum coordinates defined as:
\begin{align}
  \vv{S} & = \vv{p}_1 + \vv{p}_2 & \vv{p}_1 & = \frac{1}{2}(\vv{S}+ \vv{D})   \\
  \vv{D} & = \vv{p}_1 - \vv{p}_2 & \vv{p}_2 & = \frac{1}{2}(\vv{S} - \vv{D}).
\end{align}
The same scheme is used for the primed variables.

Here are then some useful identities that will be used later on. Firstly, the sum of the square of momentas in each axes:
\begin{align}
  p_{1x}^2 + p_{2x}^2
   & = \ab(\frac{1}{2}(S_x + D_x))^2 + \ab(\frac{1}{2}(S_x - D_x))^2         \\
   & = \frac{1}{4}\ab(S_x^2 + 2 S_x D_x + D_x^2 + S_x^2 - 2 S_x D_x + D_x^2) \\
   & = \frac{1}{4}(2 S_x^2 + 2 D_x^2) = \frac{1}{2}(S_x^2 + D_x^2)
\end{align}
In a similar manner,
\begin{equation}
  p_{1y}^2 + p_{2y}^2 = \frac{1}{2}(S_y^2 + D_y^2).
\end{equation}
Secondly, the difference of momentas squared:
\begin{align}
  p_{1x}^2 - p_{2x}^2
   & = \ab(\frac{1}{2}(S_x + D_x))^2 - \ab(\frac{1}{2}(S_x - D_x))^2      \\
   & = \frac{1}{4}(S_x^2 + 2 S_x D_x + D_x^2 - S_x^2 + 2 S_x D_x - D_x^2) \\
   & = \frac{1}{4}(4 S_x D_x) = S_x D_x.
\end{align}
Similarly,
\begin{equation}
  p_{2x}^2 - p_{2y}^2 = S_y D_y.
\end{equation}
Hence,
\begin{align}
  \epsilon(\vv{p}_1) + \epsilon(\vv{p}_2)
   & = \frac{\hbar^2}{2m}(p_{1x}^2 - p_{1y}^2 + p_{2x}^2 - p_{2y}^2)        \\
   & = \frac{\hbar^2}{2m}\ab((p_{1x}^2 + p_{2x}^2) - (p_{1y}^2 + p_{2y}^2)) \\
   & = \frac{\hbar^2}{4m}(S_x^2 + D_x^2 - S_y^2 - D_y^2).
\end{align}
By detailed balance, \(\epsilon(\vv{p}_1') + \epsilon(\vv{p}_2')\) is also equal to this. For the difference in energy,
\begin{align}
  \epsilon(\vv{p_1}) - \epsilon(\vv{p}_2)
   & = \frac{\hbar^2}{2m}(p_{1x}^2 - p_{1y}^2 - (p_{2x}^2 - p_{2y}^2)) \\
   & = \frac{\hbar^2}{2m}(p_{1x}^2 - p_{2x}^2 - (p_{1y}^2 - p_{2y}^2)) \\
   & = \frac{\hbar^2}{2m}(S_x D_x - S_y D_y).
\end{align}
Similarly,
\begin{equation}
  \epsilon(\vv{p}_1') - \epsilon(\vv{p}_2') = \frac{h^2}{2m}(S_x' D_x' - S_y' D_y').
\end{equation}

For ease of reading, we shall restate the integral:
\begin{align}
  \mathcal{J}[f]
   & = \begin{multlined}[t][0.8\textwidth]
         - \frac{1}{(2\cpi)^6} \idotsint \odif{\vv{p}_2,\vv{p}_1',\vv{p}_2'} \times \frac{(2\cpi)^3}{\hbar} |V|^2 \\
         \times \ddel(\vv{p}_1 + \vv{p}_2 - \vv{p}_1' - \vv{p}_2')
         \ddel(\epsilon(\vv{p}_1) + \epsilon(\vv{p}_2) - \epsilon(\vv{p}_1') - \epsilon(\vv{p}_2')) \\
         \times F_{121'2'}\ab\big[\psi(\vv{p}_1) + \psi(\vv{p}_2) - \psi(\vv{p}_1') - \psi(\vv{p}_2')].
       \end{multlined} \\
   & = \begin{multlined}[t][0.8\textwidth]
         - \frac{|V|^2}{8 \cpi^3 \hbar} \idotsint \odif{\vv{p}_2,\vv{p}_1',\vv{p}_2'} \\
         \times \ddel(\vv{p}_1 + \vv{p}_2 - \vv{p}_1' - \vv{p}_2')
         \ddel(\epsilon(\vv{p}_1) + \epsilon(\vv{p}_2) - \epsilon(\vv{p}_1') - \epsilon(\vv{p}_2')) \\
         \times F_{121'2'}\ab\big[\psi(\vv{p}_1) + \psi(\vv{p}_2) - \psi(\vv{p}_1') - \psi(\vv{p}_2')].
       \end{multlined}
\end{align}
Performing the change of coordinates, we get
\begin{align}
  \mathcal{J}[f]
   & = \begin{multlined}[t][0.8\textwidth]
         - \frac{|V|^2}{8 \cpi^3 \hbar} \idotsint \odif{\vv{p}_2,\vv{p}_1',\vv{p}_2'} \times \delta(\vv{S} - \vv{S}') \\
         \times \frac{4m}{\hbar^2} \delta(S_x^2 + D_x^2 - S_y^2 - D_y^2 - S_x'^2 - D_x'^2 + S_y'^2 + D_y'^2) \\
         \times F_{121'2'}\ab\big[\psi(\vv{p}_1) + \psi(\vv{p}_2) - \psi(\vv{p}_1') - \psi(\vv{p}_2')].
       \end{multlined}    \\
   & = \begin{multlined}[t][0.8\textwidth]
         - \frac{m|V|^2}{2 \cpi^3 \hbar^3} \idotsint \odif{\vv{p}_2,\vv{p}_1',\vv{p}_2'} \times \delta(\vv{S} - \vv{S}') \\
         \times \delta(S_x^2 + D_x^2 - S_y^2 - D_y^2 - S_x'^2 - D_x'^2 + S_y'^2 + D_y'^2) \\
         \times F_{121'2'}\ab\big[\psi(\vv{p}_1) + \psi(\vv{p}_2) - \psi(\vv{p}_1') - \psi(\vv{p}_2')].
       \end{multlined}
\end{align}

We then change the measure from \((\vv{p}_1', \vv{p}_2')\) to \((\vv{S}', \vv{D}')\). The jacobian can be calculated as follows:
\begin{equation}
  \begin{vmatrix}
    \pdv{S_x'}{p_{1x}'} & \pdv{S_x'}{p_{1y}'} & \pdv{S_x'}{p_{2x}'} & \pdv{S_x'}{p_{2y}'} \\
    \pdv{S_y'}{p_{1x}'} & \pdv{S_y'}{p_{1y}'} & \pdv{S_y'}{p_{2x}'} & \pdv{S_y'}{p_{2y}'} \\
    \pdv{D_x'}{p_{1x}'} & \pdv{D_x'}{p_{1y}'} & \pdv{D_x'}{p_{2x}'} & \pdv{D_x'}{p_{2y}'} \\
    \pdv{D_y'}{p_{1x}'} & \pdv{D_y'}{p_{1y}'} & \pdv{D_y'}{p_{2x}'} & \pdv{D_y'}{p_{2y}'}
  \end{vmatrix} = \begin{vmatrix}
    1 & 0 & 1  & 0  \\
    0 & 1 & 0  & 1  \\
    1 & 0 & -1 & 0  \\
    0 & 1 & 0  & -1
  \end{vmatrix} = 4.
\end{equation}
Hence,
\begin{equation}
  \odif{\vv{p}_1',\vv{p}_2'} = \ab|\jdv{p_{1x}',p_{1y}',p_{2x}',p_{2y}'}{S_x',S_y',D_x',D_y'}|\odif{\vv{S},\vv{D}} = \frac{1}{4}\odif{\vv{S},\vv{D}},
\end{equation}
\begin{align}
  \mathcal{J}[f]
   & = \begin{multlined}[t][0.7\textwidth]
         \frac{m |V|^2}{2 \cpi^3 \hbar^3} \idotsint \odif{\vv{p}_2,\vv{S'},\vv{D}'} \times \\
         \frac{1}{4} \delta(\vv{S} - \vv{S}')\delta(S_x^2 + D_x^2 - S_y^2 - D_y^2 - S_x'^2 - D_x'^2 + S_y'^2 + D_y'^2) \\
         \times F_{121'2}\ab\big[\psi(\vv{p}_1) + \psi(\vv{p}_2) - \psi(\vv{p}_1') - \psi(\vv{p}_2')]
       \end{multlined}                       \\
   & = \begin{multlined}[t][0.7\textwidth]
         \frac{m|V|^2}{8 \cpi^3 \hbar^3} \idotsint \odif{\vv{p}_2,\vv{D}'} \times \delta\ab((D_x^2 - D_x'^2) - (D_y^2 - D_y'^2)) \\
         \times F_{121'2}\ab\big[\psi(\vv{p}_1) + \psi(\vv{p}_2) - \psi(\vv{p}_1') - \psi(\vv{p}_2')].
       \end{multlined}.
\end{align}

Under this coordinate, the Fermi factor becomes
\begin{equation}
  \frac{1}{4\ab[\cosh\xi + \cosh\ab(\frac{\beta\hbar^2}{4m}(S_x D_x - S_y D_y))]
    \ab[\cosh\xi + \cosh\ab(\frac{\beta\hbar^2}{4m}(S_x' D_x' - S_y' D_y'))]}
\end{equation}
where \(\xi\) now becomes \(\frac{\beta\hbar^2}{8m}(S_x^2 + D_x^2 - S_y^2 - D_y^2) - \beta\mu\).
